\documentclass[11pt]{article}

\usepackage[utf8]{inputenc}
\usepackage[T1]{fontenc}
\usepackage[margin=1in]{geometry}
\usepackage{xcolor}
\usepackage{enumitem}
\usepackage{hyperref}
\usepackage{parskip}
\usepackage{mdframed}

% ----- Colors -----
\definecolor{revcolor}{RGB}{0,0,0}        % reviewer text: black
\definecolor{replycolor}{RGB}{0,90,180}   % author reply: blue

\hypersetup{
    colorlinks=true,
    linkcolor=blue,
    urlcolor=blue,
    citecolor=blue
}

% ----- Reviewer comment block -----
\newcommand{\reviewer}[1]{\par\medskip\noindent\textcolor{revcolor}{#1}\par}

% ----- Author reply block (colored, boxed) -----
\newenvironment{reply}
{\par\smallskip\begin{mdframed}[linecolor=replycolor,linewidth=0.6pt,
   backgroundcolor=replycolor!5,leftmargin=0pt,rightmargin=0pt,
   innertopmargin=6pt,innerbottommargin=6pt]%
   \color{replycolor}\textbf{Response:}\ }
{\end{mdframed}\par\medskip}

\title{Response to Reviewers}
\author{Maioli et al.}
\date{}

\begin{document}

\maketitle

\noindent We thank the reviewers for their careful and constructive comments.
Below, reviewer comments are shown in \textcolor{revcolor}{black} and our
responses in \textcolor{replycolor}{blue}.

%==================================================================
\section*{Reviewer \#1}
%==================================================================

\subsection*{AI Disclosure}
\reviewer{I did not use AI software during preparation of this review.}

\subsection*{General Comments}

\reviewer{In the manuscript by Maioli et al. presents a new analysis of previously
compiled and analyzed bottom trawl data from the northern Atlantic and Pacific
oceans to study evidence for spatial range shifts in marine species. The article
is well written and the analysis is appropriately sophisticated and accounts for
many of the problems inherent in complex spatio-temporal and data integration
analyses.}

\reviewer{In terms of novelty, the analysis leans into recent calls (e.g.,
Fredston et al. 2025, \textit{Trends in Ecology and Evolution}) to examine range
shifts across multiple dimensions and follows along similar analyses previously
done over large spatial scales with terrestrial organisms (e.g., Neate-Clegg et
al. 2024, \textit{Nature Ecology and Evolution} [uncited]). The present manuscript
is thus laudable for considering latitude and depth shifts simultaneously, yet
still fails to account for many other dimensions or ways that species can track
thermal niches, including aquatic microclimates (Lawlor et al. 2024, \textit{Nature
Reviews Earth \& Environment}) and phenology (Fredston et al. 2025; Neate-Clegg et
al. 2024). As other authors have shown, species are tracking thermal niches
simultaneously in many dimensions, and while the present manuscript accounts for 2
major axes, it is far from exhaustive or complete. Thus, all discussions of
tracking (or not) of thermal niches must include this rather considerable caveat.}

\begin{reply}
We included this point in the introduction.
\end{reply}

\reviewer{Relatedly, my biggest concern with this manuscript is the way that
population-level trends are aggregated to larger scales. My understanding is that
this is done hierarchically, such that you have, for example, a fish species
represented in 5 regions, each with its own variety of spatial shifts, and then a
hierarchical component looks to see if there is an ``average'' shift happening
across all those regions. While this hierarchical method makes sense in many
contexts, it can be problematic in the context of measuring range shifts, as
species are not fully contained within single regions, such that single-region
analyses of ``range shifts'' can be highly biased by both false positives and
negatives. This problem (which is pervasive within the range shift literature) is
of growing concern and recognition, and is only discussed partly by Lawlor et al.
(2024) and in this preprint by Fredston
(\url{https://ecoevorxiv.org/repository/view/7383/}). While region-specific
heterogeneity may be ``truth'', what would make me more confident in the results
would be if a secondary analysis were done, where species (across all sampled
areas) are the unit of analysis, and all data for each species are aggregated into
a single range shift metric per species. While such an analysis would still have
many issues -- e.g., measuring range shifts from discontinuous spatial data (Li et
al. 2025, \textit{Plants} [14:2]) -- at least there would be less bias from what we
expect to be considerable region-specific `noise'.}

\begin{reply}
it is  part of the paper is to find regional fingerprints which couldn't be reached without separating into region. Second survey integration is a real deal. different surveys different countries, they sample in different time gear have different catchabilities. While we put some effort in aggrageting survey when data allowed (e.g. similar survey protocolos, overlapping areas) combining across bigger area would encounter several problems like space x catchability confounding. We think one the values of this paper is precisely that, combine what we could combine rather than aggragating disparate data. We discuss that in 565 620. 
\end{reply}

\reviewer{Finally, I worry the framing of the conclusions relative to the thermal
niche (e.g., line 419--425). The study finds that most species are not `keeping up'
perfectly with shifting isotherms (a result that has been demonstrated broadly
before many times, e.g. Chen et al. 2011 \textit{Science}). The authors then
conclude that ``realized thermal niches are dynamic and adaptive under climate
change''. I support the term ``dynamic'' but it remains unclear what the
consequences are of ``dynamic realized thermal niches.'' There is considerable
debate about what the realized thermal niche even means, particularly when applied
simply to populations and not full species. There is no evidence (here, at least),
that a dynamic realized niche is ``adaptive'', particularly as the study only
monitors occurrence of species, not demography. This section should be revised.}

\begin{reply}
foloowing yours and rev 2 suggestions we avoid the word thermal niche and simply refer to the mean temperature occupied by a species.
\end{reply}

\subsection*{Line Edits}

\reviewer{\textbf{133.} You need more explanation of why horizontal (longitudinal?)
shifts may expected. Presumably that would be shifts from off-shore to near-shore,
which could impact things like SST, but is ``horizontal'' the appropriate
dimension? For example, if it's the western Atlantic versus eastern Pacific,
offshore-to-nearshore is either E-W or W-E.}

\begin{reply}
horizontal are latitudinal and longitudinal (see below). as we have substatntial restructing of the intoduciton this part is no more present
\end{reply}

\reviewer{\textbf{142.} Ah, perhaps you're defining ``horizontal'' as either
longitudinal or latitudinal? I'm not sure that's the best term -- at least I'm
finding it confusing (especially since this single dimension is then analyzed later
as two dimensions, lat and lon).}

\begin{reply}
Yes, we used horizontal as latitude and longitude following e.g. Gruenburg et al. 2025
\textit{Nature Climate Change}, we now removed many of this wording and we spell out lat and lon. where it remained we mate that explicit line
\end{reply}

\reviewer{\textbf{145.} I dislike claims of `novel' without first doing an
extensive literature review. Yes, `multi-dimensional response' papers are rare, but
they are not absent from the literature (e.g., Neate-Clegg et al. 2025). The impact
of this work is not lessened by simply saying, ``Unlike some previous
efforts\ldots, our approach allows\ldots''}

\begin{reply}
we removed novel from the sentence
\end{reply}

\reviewer{\textbf{679.} Your methods imply here, if I'm following them correctly,
that you estimated range shifts not from data or models directly, but by projecting
data onto space, and quantifying range shifts that way? If so, doesn't that mean
that your range shifts are potentially based substantially on extrapolation (or
interpolation), particularly to areas with unsampled environmental variable
combinations? In other words, it seems your method is akin to measuring range
shifts from projecting (complex spatio-temporal) SDMs across space and time and
measuring range shifts from those projections? Such a method seems very
susceptible to false positives (which I strongly feel has been shown before in the
literature, but not sure of the reference).}

\begin{reply}
We do compute distributional metrics from model-based predictions, but this does not involve extrapolation. Our prediction grid is restricted to the surveyed footprint of each region (both space and time), and grid-cell covariates span only the range of depths and temperatures observed there. We predict within the sampled domain, not into novel environmental space.

Model-based estimation tends to reduce rather than introduce this bias. Raw metrics are sensitive to changes in the spatial distribution of sampling effort over time, which can produce apparent centroid shifts in the absence of real redistribution; estimating the underlying density field first accounts for this and yields less biased centre-of-gravity estimates than design-based alternatives (Thorson et al. 2016, ref 40).
\end{reply}

%==================================================================
\section*{Reviewer \#2 (Alexa Fredston)}
%==================================================================

\subsection*{Comments on Significance Statement}
\reviewer{This is one of my favorite parts of the draft. It is well written and
accurate.}

\subsection*{Comments}

\reviewer{This paper revisits several core findings in marine global change biology
in the past 20 years---specifically, that species are deepening and shifting
geographically in response to climate change. They bring new methods to bear on
this question, and in particular, substantial advances in both how the survey data
used to answer this question are treated as well as statistical modeling that
parses out species- and regional-level variation from temporal trends much better
than past attempts.}

\reviewer{Before I get into the paper I want to commend the authors on the
accessibility and documentation of their data and code. In papers that hinge on
statistical modeling it's critical for reviewers to actually read the code and
understand its concordance with the paper, and too often I see authors obscure
their data and code to prevent exactly this level of scrutiny. This clearly took a
lot of work to organize so nicely on GitHub and it should not go unrecognized.}

\reviewer{Here's my summary (and correct me if I'm wrong) of the methods and
results:}

\reviewer{%
\begin{enumerate}
    \item The authors started with a large database synthesizing bottom-trawl
    survey records of marine fishes, pooling surveys where they could and
    accounting for variation between them when they couldn't.
    \item They then fitted species distribution models (one species at a time) to
    turn these observation points into annual estimates of biomass density over
    space, with associated uncertainty.
    \item In a hierarchical framework that propagated this uncertainty and also
    allowed for correlations both among species within a region and among response
    variables for a given species, they tested whether four metrics derived from
    the SDMs were changing over time:
    \begin{enumerate}
        \item Latitude of centroid (center of gravity, AKA weighted mean, of the SDM)
        \item Longitude of centroid
        \item Mean biomass-weighted depth
        \item Biomass-weighted mean of annual mean sea bottom temperature occupied
    \end{enumerate}
    \item They find (with a colossal dataset and decades of data) that none of the
    first three response variables are changing significantly over time at the
    global (well, Northeast Pacific / North Atlantic) scale; and most of them
    aren't even changing significantly over time at the regional scale.
    \item This result implies that fish species are on average (but not necessarily
    individually) sitting in place while the oceans warm, which means the mean
    temperature they occupy is getting hotter.
\end{enumerate}}

\begin{reply}

\end{reply}

\reviewer{Hopefully, readers of this manuscript will appreciate that while
calculating metrics of change without first modeling spatiotemporal patterns in the
data sometimes leads to greater statistical significance, that approach (which I am
guilty of, earlier in my career!) also makes it impossible to disentangle
observation and process error. Additionally, because all the response variables are
averages, they are by definition quite insensitive to outliers. In other words,
using SDMs and mean response variables probably made it harder to find big
significant results here, which is a feature---not a bug---of this approach. The
authors set up this analysis to test ``if we control for everything else, do we
still see big range or depth shifts in the biggest-ever dataset used to test
this?'' and the answer is no.}

\reviewer{I have some minor notes on methods and documentation, and some line
edits, below. But at a high level, I have no issues with the modeling approach used
here. I do have some issues with the framing of the text, so let me start there.}

\reviewer{The Introduction makes it sound like the authors are tackling the very
important and emergent question of simultaneously modeling species' responses to
environmental change along multiple possible gradients to fully quantify their
degree of thermal niche tracking (sensu Neate-Clegg et al. 2024 \textit{Nature
Ecology and Evolution} and Fredston et al. 2025 \textit{TREE}). But this isn't what
the paper does at all. The models in this paper can't quantify the degree to which
shifts along one dimension (depth vs longitude vs latitude) are trading off against
one another to achieve climate tracking in sum. Which is fine---this paper does
something else, which is equally cool and just different---but I think the
Introduction and really the whole paper need substantial reframing.}

\reviewer{For example, the second paragraph of the Introduction (lines 51--61) set
up the knowledge gap this paper will fill as [to paraphrase] ``perhaps species
aren't uniformly marching towards the poles!'' This is a strawman argument that is
over a decade out of date, and it doesn't do the models in this paper justice. The
models in the paper also can't really quantify the degree to which fish are
responding coherently to warming (L58--61), because they aren't actually
quantifying thermal niche tracking over time.}

\begin{reply}
The reviewer is right that we ran two claims together, and we have separated them.
One is whether species move coherently — whether species in a region share a direction, and whether regional patterns add up across basins. Our model does estimate this. It splits each trend into a global effect, a regional deviation and a species deviation, so how much of the variation is shared is itself a measure of coherence. Figure 3 shows the same thing more directly, as the proportion of species moving the same way. The Introduction now attaches this claim to the model, rather than thermal tracking.
The other is tracking. A species keeping up with its isotherm would stay at the same temperature; the fishes here did not, and the temperatures they occupy rose (well, 60\% of them). That is a measure of tracking, in degrees rather than kilometres. Comparing it with a static distribution then shows how much of the rise came from the water warming rather than from where the fish were.
What we cannot do is say how any compensation was divided between latitude, longitude and depth, or judge movement against the distance an isotherm actually travelled. We have cut the sentences that implied otherwise. The Introduction now says only that we estimate whether movement along the different axes is coordinated, and what it did to the temperatures species occupy.
\end{reply}

\reviewer{To expand on that point, I was puzzled by the use of ``niche'' to
describe a scalar mean of means. The niche, by definition, is something with a
spread or a breadth; niche theory has never suggested that species occupy only a
single degree-Celsius temperature. Thermal niches in global change biology are
usually defined by their endpoints and often studied in terms of whether past,
present, and future habitats are inside or outside of the niche (some relevant
papers: La Sorte and Jetz 2012 \textit{Journal of Animal Ecology}; Early and Sax
2014 \textit{Global Ecology and Biogeography}; Godsoe 2010 \textit{Oikos};
Chevalier et al. 2024 \textit{Nature Ecology and Evolution}; Socolar et al. 2017
\textit{PNAS}; this is a random sample and there is no need to cite any/all of them,
but I'm including them to underscore that this interpretation of what ``niche''
means is not just my own). This paper does not need to be about niches (and in my
opinion, it isn't). I would edit out the word and just replace it with ``Depth
centroid'' and ``Mean temperature occupied'' or whatever more descriptive terms the
authors choose. I would also change the title, which fails to capture the core
results of the paper. (Actually, the only part of the paper which I think is worded
in a way that reflects the analysis is the significance statement, which is great.)}

\begin{reply}
We agree. A niche has breadth, and what we estimate is a single density-weighted mean of the temperatures a population occupies — one summary of that distribution, not the distribution itself. Calling it a niche claimed more than the analysis supports.

We have removed the word from the paper. "Depth niche" is now "depth centroid" and "thermal niche" is now "mean occupied bottom temperature", following the reviewer's suggestion. The change runs through the title, abstract, main text, Table 1, all figure labels and the supplement.

We have also retitled the paper "[NEW TITLE]", which states what we found rather than the framework we used. We were glad the significance statement reads well, and have used it as the guide for wording the rest.
\end{reply}

\reviewer{So, how should these results be framed? I would present this paper by
saying that past attempts to quantify the true degree to which species are shifting
latitudinally, longitudinally, and in depth were limited---in their geographical
scope, their data volume, and their simplistic statistical approaches. So we are
lacking strong tests of the hypotheses about thermal tracking in marine systems set
up by Sunday et al. 2012 \textit{Nature Climate Change}, Burrows et al. 2011
\textit{Science}, and others. Further, estimated rates of shifts usually come from
meta-analyses of published studies (Lenoir et al. 2020, Poloczanska et al.
2013)---not from directly analyzing synthesized datasets---and those are rife with
confounding variables and methodological issues (see Brown et al. 2016
\textit{Global Change Biology}, from the same team as Poloczanska et al., or Lawlor
et al. 2024, from the same group as Lenoir et al.). No one has recently brought
updated data or statistics to the question ``How fast are marine fishes actually
shifting or deepening?'' and in my opinion that is the great contribution of this
paper, which finds that actually they aren't consistently doing either (although
individual species often move a lot), and because the oceans are warming so too are
the temperatures where these fishes hang out.}

\reviewer{As a proud author of two FISHGLOB-based studies with null results that
also disproved beautiful hypotheses (Fredston et al. 2023 \textit{Nature} and
Kitchel et al. 2025 \textit{PLOS Climate}), I understand that it's hard to write
global change papers around the finding that nothing changed. But my opinion is that
it's absolutely critical to write and publish these papers, particularly in an
instance with such careful modeling (and I give PNAS kudos for sending this out for
review). Of course, the authors probably can imagine other framings that would suit
this paper too, but my key point is that the current framing both sells the actual
results short and is a bit misleading.}

\reviewer{The last issue I will raise is the choice to frame this around climate
change responses when many of the focal regions did not warm in the study interval.
It would be strange if the authors had found consistent directional shifts in half
of these regions, where mean sea bottom temperatures were relatively constant (see
Fig. 1). There is no mention of these regional differences in warming in the paper,
which likely have a huge impact on the results. The authors might want to, for
example, report and interpret results separately by whether the region cooled,
warmed, or neither.}

\begin{reply}
We have made the regional differences in warming explicit. Figure 1 now reports each region's bottom-temperature trend with a credible interval, and the text states which regions warmed and by how much.

We would also argue the variation is useful rather than a problem. Regions that warmed little provide the contrast that lets us ask whether responses scale with warming — and they do. Trends in occupied temperature follow local warming rates closely across regions ($\rho$ = 0.92), and the three regions with no detectable warming of occupied temperature are among those where bottom temperatures changed least. The paper does not treat all regions as though they experienced the same climate signal.

Following the reviewer's suggestion, we now report results separately for regions that warmed over the study period and those that did not. [RESULT.] We have added this to the Results and revised the Discussion accordingly.
\end{reply}

\subsection*{Line and minor edits, in no particular order}

\reviewer{The Discussion should probably address why the authors didn't quantify
concordance of biotic shifts with climate velocities, an approach that was invented
to deal with regionally-divergent patterns in redistribution (Pinsky et al. 2013
\textit{Science}).}

\begin{reply}

\end{reply}

\reviewer{The Discussion (or Introduction) should also be sure to differentiate
this paper from others using the FISHGLOB dataset, notably Gruenburg et al. 2025
\textit{Nature Climate Change}.}

\begin{reply}

\end{reply}

\reviewer{Nowhere in the paper does it say the actual year range of the data, which
should be included!}

\begin{reply}

\end{reply}

\reviewer{Fig. 1 should show the significance of the plus/minus temperature change
rates printed in the panels, and/or report a credible interval, since I suspect
some (many?) overlap with zero.}

\begin{reply}

\end{reply}

\reviewer{In a sentence about marine species tracking temperature change, why cite a
paper about birds (ref 6)?}

\begin{reply}

\end{reply}

\reviewer{It's not clear what ``net coherent redistribution signal'' means; I would
either define this term, or edit it out and replace it with a descriptive sentence.}

\begin{reply}

\end{reply}

\reviewer{A rationale should be included for the 15-year cutoff for dataset length.}

\begin{reply}

\end{reply}

\reviewer{In the equations and text around them, ``observation $i$'' is not defined,
which makes it hard to understand the equations.}

\begin{reply}

\end{reply}

\reviewer{I would appreciate a table in the SI articulating the correspondence
between terms in equations in the text and named objects in the code. I know it's
impossible to fully describe models in text, and I often go look at the code to
figure out what really happened, but here it was hard to map them onto one
another---for example, is \texttt{decade\_i} in the paper the same as
\texttt{years\_c} in the code (or \texttt{years\_c} times ten, or something else)?
Table 1 in this Supplement is an example of what I mean:
\url{https://github.com/afredston/mid_atlantic_forecasts/blob/main/09_supplement.pdf}.}

\begin{reply}

\end{reply}

\end{document}
